% Solutions to Chapter 11's exercises, from lectures/L11/appendix/d_solutions.md.
%
% The programs are left out: exercise 7 c) keeps its answer to the question whether r24 can be
% loaded once, and 7 d) says in words how the program it counts is built; exercise 8 keeps the flow
% chart in words, what the simulator showed, and c). Exercise 9 c) said the difference was "ett
% tiotal procent"; with the solution's own figures, 12,7 mA against 13,6 mA, it is about 7 %, and it
% now says so, here and in the lecture. Exercise 3 b) called an 18 mA margin "knappt en lysdiod
% till", less than one more LED; one LED takes 13,6 mA, so it is "drygt en lysdiod till", here
% and in the lecture.

\answerchapter{c11:ch}

De cykeltal som anges här är uppmätta i simulatorn.

\answer{c11:ex:1}{Ingångar, beslut och utgångar}
\pt{a}Ingångar: knappen och de två gränslägesbrytarna. Utgångar: motor upp och motor ned. Beslut:
ett tryck startar porten mot det läge den inte står i, och motorn stannar när gränslägesbrytaren i
det läget slår till. Ett bra program ser också till att upp och ned aldrig är på samtidigt.

\pt{b}Ingång: IR-givaren. Utgång: fläkten. Beslut: fläkten går så länge givaren ser en hand, och
tio sekunder till efter att den försvunnit. Beslutet innehåller alltså en tid, precis som
trafikljuset.

\pt{c}Ingång: knappen. Utgångar: den röda, den gula och den gröna lampan. Beslut: grönt tills
knappen trycks, sedan en fast följd av tillstånd med bestämda tider, och grönt igen.

\answer{c11:ex:2}{Förkopplingsmotståndet}
\pt{a}Spänningen över motståndet är \code{5 - 2 = 3 V}.
\begin{itemize}
  \item 220~Ω: \code{3 / 220 ≈ 0,0136 A = 13,6 mA}
  \item 330~Ω: \code{3 / 330 ≈ 9,1 mA}
  \item 1~kΩ: \code{3 / 1000 = 3 mA}
\end{itemize}

\pt{b}\code{(5 - 3) / 220 ≈ 9,1 mA}. Den gröna lysdioden tar mer av spänningen själv, så det blir
mindre kvar över motståndet och därmed mindre ström. Dessutom ger olika färger olika mycket ljus
för samma ström, så gröna lysdioder upplevs ofta som svagare.

\pt{c}\code{R = 3 V / 0,010 A = 300 Ω}. Närmaste vanliga värde är 330~Ω, som ger 9,1~mA, eller
270~Ω, som ger 11~mA. Båda duger.

\pt{d}\code{3 / 47 ≈ 64 mA}, långt över 40~mA, som är stiftets absoluta maximum. Både lysdioden
och stiftet riskerar att förstöras. I praktiken sjunker spänningen på stiftet vid så stor ström, så
det blir något mindre, men det är ändå långt utanför vad kretsen är specificerad för.

\answer{c11:ex:3}{Hela porten}
\pt{a}\code{6 · 13,6 mA ≈ 82 mA}.

\pt{b}Marginalen är ungefär 18~mA, drygt en lysdiod till. Det fungerar, men det är inte mycket att
ha om någon senare lägger till fler utgångar på samma grupp.

\pt{c}Med 330~Ω blir det \code{6 · 9,1 ≈ 55 mA}, med god marginal, och lysdioderna lyser nästan
lika starkt. Att välja det lite större motståndet kostar nästan ingenting och tar bort risken.

\answer{c11:ex:4}{Reläet}
\pt{a}En reläspole behöver typiskt mer ström än ett stift kan ge, ofta 70~mA eller mer för ett relä
på 5~V. Transistorn låter några milliampere från stiftet styra den större strömmen.

\pt{b}När transistorn slår av vill spolen hålla sin ström igång, och spänningen över den slår upp
till många tiotals volt. Dioden leder då och låter strömmen klinga av genom spolen. Utan dioden går
spänningstoppen över transistorn, som kan förstöras redan första gången.

\pt{c}Störningar och fel på 24~V-sidan, en kortslutning, en spänningstopp från en motor, kan inte
nå mikrodatorn. De två kretsarna kan dessutom ha olika jord.

\pt{d}Stiftet ger bara 5~V, så lampan kommer knappt att glöda. Men en kall glödtråd har lågt
motstånd, och strömmen kan bli långt över vad stiftet tål. Stiftet, och i värsta fall kretsen,
riskerar att förstöras. Rätt sätt är ett relä, som i \secref{c11:a3}.

\answer{c11:ex:5}{Knappar som studsar}
\pt{a}Programmet räknar varje gång stiftet går från 1 till 0. I testet slöts knappen två gånger vid
varje tryck, sluts, studs, sluts igen, och en gång till när den släpptes. Varje sådan nolla ser för
programmet ut som ett nytt tryck: tre per tryck, sex för två tryck.

\pt{b}Väntan ska vara längre än studsarna varar och kortare än den tid en människa behöver mellan
två tryck.
\begin{itemize}
  \item \textbf{1~ms} är kortare än studsarna, som kan hålla på i flera millisekunder. Programmet
    tittar på knappen igen medan den fortfarande studsar, och räknar för mycket.
  \item \textbf{500~ms} tar bort studsarna, men programmet blir trögt: ett andra tryck inom en halv
    sekund missas, och det märks för den som trycker snabbt.
\end{itemize}

\pt{c}Programmet räknar ingenting. Under några millisekunder av studs flimrar lysdioden, vilket
ingen människa ser, och därefter följer den knappen.

\answer{c11:ex:6}{Hur lång tid tar ett varv?}
\pt{a}\mbox{\code{16 000 084 · 4 + 18 = 64 000 354}} cykler.

\pt{b}Ett tillstånd tar \mbox{\code{1 + 1 + 1 + 64 000 354 = 64 000 357}} cykler. Två tillstånd
och ett \code{rjmp}: \mbox{\code{2 · 64 000 357 + 2 = 128 000 716}} cykler. Precis vad kontrollen
mätte.

\pt{c}8~s är 128\,000\,000 cykler, så varvet är 716 cykler för långt. En cykel är 62,5~ns, så
\mbox{\code{716 · 62,5 ns ≈ 44,8 µs}}. En timme är \mbox{\code{3600 / 8 = 450}} varv, och
\mbox{\code{450 · 44,8 µs ≈ 20 ms}}.

\pt{d}Resonatorn kan ge 0,5~\%, alltså 18 sekunder på en timme. Det är nästan tusen gånger mer än
programmets 20~ms. Det är resonatorn som är värd att bry sig om, om noggrannheten alls spelar roll,
och för ett trafikljus gör den inte det.

\answer{c11:ex:7}{Nattläge}
\pt{a}\leavevmode

\begin{narrowtable}{@{}lccl@{}}
  \thead{Tillstånd} & \thead{Gul (PB1)} & \thead{Tid} & \thead{Därefter} \\
  \midrule
  Tänd & tänd & 0,5~s & Släckt \\
  Släckt & släckt & 0,5~s & Tänd \\
\end{narrowtable}

\pt{b}Start, initiering av \code{SP} och PB1, och sedan en loop utan villkor: tänd gul, vänta
0,5~s, släck gul, vänta 0,5~s, tillbaka till \enquote{tänd gul}. Samma form som vägarbetsljusets
flödesplan.

\pt{c}Ja, \code{r24} kan laddas en enda gång, eftersom \code{delay_ms} följer kursens regel och
lämnar \code{r24} oförändrat. Med en subrutin som inte gjorde det hade \code{r24} behövt laddas om
före varje anrop.

\pt{d}Räkna med ett program där \code{r24} = 250 laddas före loopen, och loopen är \code{sbi},
två anrop av \code{delay_ms}, \code{cbi}, två anrop till och ett \code{rjmp} tillbaka. Den tända
halvan är \code{sbi} (2 cykler) och två anrop om \mbox{\code{4 000 018}} cykler:
\mbox{\code{8 000 038}}. Den släckta är \code{cbi} (2), två anrop och \code{rjmp} (2):
\mbox{\code{8 000 040}}. En hel period är \textbf{16\,000\,078} cykler, en sekund och knappt
5~µs. Simulatorn mäter samma sak.

\answer{c11:ex:8}{En knapp som växlar}
\pt{a}Start och initiering, sedan en loop: vänta tills knappen är nedtryckt, växla lysdioden, vänta
20~ms, vänta tills knappen är släppt, vänta 20~ms, och tillbaka.

\pt{b}I simulatorn, med samma studsande tryck som \code{press_counter.asm}, ska lysdioden lysa
efter det första trycket och vara släckt efter det andra.

\pt{c}Utan väntan växlar lysdioden en gång för varje studs. Är antalet studsar jämnt ser det ut som
att trycket inte gjorde någonting, och är det udda fungerar det. Resultatet blir slumpmässigt, och
det är just det som gör felet så förvirrande att hitta.

\answer{c11:ex:9}{Kontroll: strömmen genom lysdioden}
\pt{a}\code{(5 - 2) / 220 ≈ 13,6 mA}.

\pt{b}Ett typiskt resultat är ungefär 2,8~V över motståndet, alltså \code{2,8 / 220 ≈ 12,7 mA}, och
ungefär 4,8~V mellan stiftet och jord. Era värden beror på lysdioden och på USB-porten.

\pt{c}Strömmen blir något mindre än uträknat, av två skäl som mätningen visar:
\begin{itemize}
  \item \textbf{Stiftet ger inte riktigt 5~V under last.} Utgången har ett litet inre motstånd, så
    spänningen sjunker när en ström går ut ur den. Dessutom är matningen från USB inte alltid
    exakt 5~V.
  \item \textbf{Lysdiodens spänning är inte exakt 2~V.} Den beror på typ och färg, och varierar
    något med strömmen.
\end{itemize}
Skillnaden är några procent, i exemplet ovan ungefär 7~\%, och den visar hur mycket en förenklad
modell räcker till: tillräckligt för att välja motstånd, inte för att förutsäga strömmen på
milliampere när.

\pt{d}Den gröna lysdioden har högre spänning över sig, kanske 2,2--3~V beroende på typ. Då blir det
mindre kvar över motståndet, och strömmen blir mindre. Mät spänningen över lysdioden, så ser ni hur
stor skillnaden är.

\answer{c11:ex:10}{Felsökning på kortet}
\pt{a}Troligast är den gröna lysdioden vänd fel, eller kopplad till fel stift. Kontrollera att det
långa benet sitter mot motståndet, och att sladden går till stift 10. Byt plats på den gröna och
den röda lysdioden: följer felet med lysdioden, är det den eller dess koppling.

\pt{b}\code{QUARTER_MS} står kvar på 1. Hela ljuscykeln går då runt ungefär 30 gånger i sekunden,
och ögat ser alla lamporna lysa svagt. Sätt den till 250, bygg om och för över programmet igen.

\pt{c}Programmet har aldrig kommit fram till kortet: fel COM-port, eller kortet är inte anslutet.
Kontrollera porten i Arduino IDE och i verktygets argument, se \secref{studio:5-6} i
bilaga~\ref{app:studio}.

\pt{d}Sladdarna till stift 8 och 9 har bytt plats, eller så har bitnumren i \code{.equ RED} och
\code{.equ YELLOW} bytt plats i programmet. Jämför kopplingen med tabellen i slutuppgiften.

\answer{c11:ex:11}{En dag på ett trafikljus}
\pt{a}\mbox{\code{0,005 · 86 400 s = 432 s}}, alltså 7 minuter och 12 sekunder på ett dygn.

\pt{b}\mbox{\code{86 400 / 8 = 10 800}} ljuscykler.
